% ===================== CHAPTER 1 — INTRODUCTION =====================
\chapter{Introduction}\label{cha:introduction}

\section{Background}

Indian Railways reserves operate through large tracts of forest and open
grassland in which scheduled and unscheduled train paths intersect with
animal movement corridors. Elephants, deer and other wildlife cross
railway formation at grade, and collisions remain a persistent operational
and ecological risk. Conventional protection --- signage, whistle zones and
manual watchmen --- is sparse by design, and a single fixed level-crossing
alarm cannot adapt on its own to \emph{which} animal is ahead, \emph{where}
it is relative to a moving train, or \emph{how much} braking distance the
train still has.

Modern practice addresses this with two converging ideas:

\begin{itemize}
  \item \textbf{Digital twins:} a live virtual replica of a physical asset
        (here, a forest railway line with its trains, signals and crossing)
        that stays synchronised with trusted state and lets operators
        observe, rehearse and reason about operations.
  \item \textbf{Embedded IoT sensing:} low-cost edge devices (here, an
        ESP32 with sensors, protective outputs and a fail-safe state
        machine) at the trackside, publishing machine observations over an
        authenticated message bus rather than treated independently.
\end{itemize}

This project, developed as an educational simulation by Team 14 for the
course \emph{Embedded Systems and Internet of Things}, brings both ideas
together in a system small enough to run on one laptop: a three-dimensional
digital twin of a forest railway loop, a trusted backend control plane, a
role-aware web dashboard, an authenticated MQTT broker, and an ESP32 edge
controller --- all exchanging sensor observations, authority commands,
state snapshots and alerts over well-defined topics.

\section{Objectives}

The objectives of the project are:

\begin{enumerate}
  \item Build a runnable \textbf{3D digital twin} of a double-track forest
        railway --- UP and DOWN loops, an island station with interlocking
        crossovers, a wildlife level crossing with boom barriers, a carved
        river with road and railway bridges, catenary, dynamic signals and
        an LTS-derived Blender asset fleet --- in Godot~4.
  \item Construct a \textbf{trusted control plane} that owns all railway
        authority: a two-train double-track fleet model, per-block signal
        aspects, a fixed-timestep simulation loop, a station-crossover
        interlocking, chance-based meet scheduling, an operator-requestable
        interlocked wrong-line run, and a wildlife-hazard state machine ---
        served over authenticated REST and MQTT.
  \item Enforce \textbf{security end to end}: authenticated users with
        role-based access control, an authenticated MQTT broker with
        per-client topic ACLs, command-envelope validation with
        timestamp freshness and anti-replay protection, and a complete
        audit trail.
  \item Implement \textbf{wildlife sensing and staged protection}: a
        NONE~$\rightarrow$~SUSPECTED~$\rightarrow$~CONFIRMED escalation
        driven by trackside sensor huts, an onboard simulated camera
        detector with analytic occlusion and a quadratic stopping-distance
        model, a backend camera ingest pipeline with two-frame evidence
        confirmation, and ML-assisted animal detection through a FastAPI
        microservice.
  \item Implement an \textbf{ESP32 edge node} with restrictive fail-safe
        outputs, debounced local inputs, SNTP-gated clock truth,
        non-blocking telemetry and edge-triggered wildlife alerts ---
        validated natively and in the Wokwi cloud simulator.
  \item Provide a \textbf{role-aware dashboard} for viewer, operator and
        admin users with live state, per-block aspects, per-train speed
        control, route/meet/wrong-line panels and hazard evidence badges.
  \item Verify the whole stack continuously with an automated
        multi-surface test gate and documented live end-to-end scenarios.
\end{enumerate}

\section{Scope and safety boundary}\label{sec:scope}

Two scope statements are as important as the features themselves:

\begin{itemize}
  \item \textbf{Educational simulation only.} The system is a classroom
        digital-twin simulation, \emph{not} railway safety equipment. It
        must never be connected to real signalling, trains or operational
        infrastructure. Simulated ``stopping distance'' is a visualization
        estimate only and does not model certified train dynamics,
        gradients, adhesion or real braking curves.
  \item \textbf{Physical hardware status.} The ESP32 path was verified
        through the Wokwi cloud simulator (via a Pinggy tunnel to the
        local authenticated broker). The physical-hardware checklist items
        in the Stage~D plan remain \emph{not run}: real instruments were
        not connected. The report is explicit wherever this applies.
\end{itemize}

Within those boundaries, every subsystem --- backend, twin, dashboard,
firmware and ML service --- is fully implemented and gated by automated
tests (Chapter~\ref{cha:features}, Appendix~\ref{cha:appendix-gates}).

\section{Phased implementation outline}

The system was built in recorded, gated stages (dates from
\texttt{docs/porting/acceptance.md}):

\begin{itemize}
  \item \textbf{Stage A (backend foundation)} --- Express + JWT/RBAC,
        SQLite audit and anti-replay registry, two-zone state machine,
        backend-authoritative MQTT, broker auth and ACLs
        (2026-10-18/19).
  \item \textbf{Stage A2 (block signalling)} --- per-block aspects and
        retained signals (2026-10-19).
  \item \textbf{Stage B (dashboard)} --- login, role-aware REST controls,
        WebSocket live updates, audit feed (2026-10-19).
  \item \textbf{Stage C (Godot twin)} --- vendored MQTT bridge, procedural
        world, train/wildlife animation, brakes, signals, cameras;
        porting Libre-TrainSim (LTS) Blender assets (2026-10-19/20).
  \item \textbf{Stage D (ESP32 / Wokwi)} --- per-client credentials,
        device-scoped topics, fail-safe outputs, cloud simulation through
        a tunnel (2026-10-19/20).
  \item \textbf{Stage E (integrated delivery)} --- Docker/environment
        setup, scenario scripts and the formal acceptance record
        (2026-10-22).
  \item \textbf{Stage 2A--2G (double-track phase)} --- a two-train UP/DOWN
        fleet, ten blocks with aspects, station island and crossovers
        (2A), the double-track world (2B--2C), backend lives and camera
        fleet (2A--2D), meet/pass scenarios (2E, 2G), and distributed
        wildlife sensing with an onboard detector (2F) (2026-10-21/22).
  \item \textbf{Stadium evolution --- Phases 1--3} --- both loops reshaped
        as one 250~m straight-sided shared stadium footprint with opposite
        operational chainage (2026-10-22), an interlocked operator
        wrong-line run with an opposing-train hold (2026-10-23), and the
        staged camera/ML/near-animal confirmation chain with two-frame
        evidence (2026-10-23).
\end{itemize}

\section{Organisation of the report}

Chapter~2 states requirements. Chapter~3 describes the system architecture
and end-to-end workflow. Chapter~4 documents tools and technology.
Chapter~5 covers hardware, peripherals and the simulation stack.
Chapter~6 details the digital twin world. Chapter~7 documents the
dashboards. Chapter~8 lists every implemented feature and function with
its verification evidence. Chapter~9 presents results, limitation and
conclusion. Appendices collect the access-control matrix, the complete
test-gate record and run instructions. Reproduced simulation screenshots
appear in Figures~\ref{fig:overview}--\ref{fig:cameraconfirm} and the
figures referenced from Chapters~5--8.
